\documentclass[12pt]{article}
\usepackage{graphicx} % Required for inserting images
\usepackage{authblk}      % Optional: better author formatting
\usepackage{setspace}     % Optional: spacing control
\usepackage{fancyhdr}
\usepackage{amsmath}
\usepackage{caption}
\usepackage{pdflscape}
\usepackage{graphicx}
\usepackage{booktabs}
\usepackage{adjustbox}
\usepackage{url}
\usepackage[utf8]{inputenc}
\usepackage{csquotes}
\usepackage{natbib}
\usepackage{pdfpages}
\usepackage[raggedright]{ragged2e}
\usepackage[hidelinks]{hyperref}
\usepackage[letterpaper,margin=1in]{geometry}
\bibliographystyle{apsr}   % or plainnat, chicago, apa, etc.
\setcitestyle{authoryear,open={(},close={)},aysep={ }}

\title{Research Design for We Don't Do That: State Identity and Nuclear Non-Use} 
\author{Wyatt King}
\date{March 16, 2026}


\begin{document}

\maketitle
\section{Research Design}
To test whether identity-based arguments and knowledge of India's NFU reduce support for nuclear use, I conduct a conjoint survey experiment modeled off of those created by Sagan and Valentino. In the experiment, respondents read a scenario describing a hypothetical Indian counterforce attack against Pakistan. The situation described is consistent with India's Cold Start doctrine and inspired by what security scholars consider to be the most likely scenario in which India would first use nuclear weapons \citep{kapur_ten_2008, clary_indias_2019}. Implemented after the Kargil War, Cold Start is military doctrine enabling India to quickly mobilize along the border with Pakistan following a state-sponsored terrorist attack. Within days of mobilization, India can push across the borders with Punjab and Kashmir either to claim disputed territory or force a favorable compromise. 

Although India publicly maintains a no first use policy, Christopher Clary and Vipin Narang argue that during such a ground war, India may contemplate launching a nuclear attack. Because India's military is conventionally superior to Pakistan's, the latter could consider using a small-yield tactical nuclear weapon on the battlefield to stop an Indian incursion. However, if Pakistan is incentivized to use nuclear weapons, India is encouraged to take first mover advantage and do the same. Fearing an eventual strike against one of its cities, Clary and Narang believe that India could launch a counter force attack, where it tries to eliminate Pakistan's strategic nuclear weapons. With the risk of a large-scale Pakistani attack marginalized, India could then conduct the ground war on its own terms.

\textit{A priori}, it is difficult to know how successful such an attack would actually be, given uncertainty over where Pakistan's nuclear weapons facilities are, how effective its air defenses would be, and the ordinary risks entailed by any large-scale military operation. In the article, I note that it is unclear whether all of Pakistan's strategic nuclear weapons have been destroyed, but that Indian government officials speaking on the condition of anonymity believe that there is a low chance of retaliation. I do so as to create a more difficult test for the norm against nuclear use while not misleading respondents into believing that a counterforce attack would be guaranteed to work. In a similar vein, although it is difficult to gauge how many civilians would die in a counterforce attack, I note that the World Health Organization warns that it could reach into the hundreds of thousands. Again, I do so as to not mislead respondents into believing that a counter force attack could be committed with little harm to civilians, even if the extent of that harm is ambiguous.

The story is structured like a mock news article so that key pieces of information can be highlighted in the headline, dek, and pull quotes.\footnote{All experimental materials can be found in  \autoref{app:vignettes}.} Specifically, I emphasize information that could influence people's understanding of the strategic utility of the attack and whether it was morally justifiable—including the target of the attack, that it was launched during an ongoing ground war, the number of civilian causalities, and the risk of retaliation. Across treatment and control groups, the article is identical, except that people assigned to learn about India's NFU read an additional sentence saying that “the nuclear attack violates long-standing Indian policy that it would never use nuclear weapons unless first attacked with weapons of mass destruction.” A pull quote in the middle of the article draws further attention to this fact.

After reading the vignette, respondents are randomly assigned to either reading or not reading three arguments against nuclear use. The arguments are written to invoke India's personal, positive identity and its negative, social identity vis-a-vis nuclear use. The first two directly mention Gandhi, Nehru, and India's legacy of non-violence. In addition to calling on these historical figures, the second and third juxtapose India with other states, saying that the use of nuclear weapons violates India's unique history as an advocate for global disarmament. The second explicitly mentions the U.S. and U.K. Mentioning the only state to have used nuclear weapons and the one that colonized the subcontinent may serve to bolster nationalist claims that India has a unique moral authority. In both the first and second argument, there are appeals in the first-person plural. Evoking a conception of “we” is important because norms are inter-subjective. It is not sufficient for one person to hold a reservation about a given course of action, but it must be held among a group of people. The third argument does not include specific mentions of Indian historical figures or “we-ness”, but refers broadly to India as having been a responsible state and upholding moral leadership. 

To measure how strongly the public supports nuclear weapons use, I include six different survey questions. In keeping with existing survey experiments, I ask respondents how strongly they approve of the nuclear attack as described in the vignette and how strongly they would prefer it to an alternative (here, continuing the ground war with Pakistan). Approval is the most direct way to measure whether someone supports nuclear use. Approval is also relevant to elected politicians fearful of domestic backlash and keen on limiting public opposition to their actions. Asking people whether they would prefer nuclear use to another military action is valuable because it reveals how strongly they have internalized just war theory. Whereas the nuclear attack described could not be defended under just war theory because of its failure to discriminate between combatants and civilians, as well as its lack of proportionality, the ground war would be defensible depending on the style and extent of combat, which is not described in the vignette. 

In addition to asking for people's approval and preferences, I ask two questions regarding what material action they would take to punish or not punish the sitting government. Specifically, would they be more or less likely to vote for the ruling party, and how likely would they be to protest nuclear use? Similar questions have yet to be asked in extant survey experiments, but offer a more direct test of whether politicians would face reprimand for nuclear use than simply asking for public approval. Though approval is a useful measure of overall public support for a politician, it is only relevant to an election-seeking politician who believes that it will affect how people vote. If a given policy decision is widely unpopular but low salience, then it will have little effect on voting patterns. By contrast, if the reaction to an action is polarized but high salience, then it could either incentivize or disincentivize a politician from supporting nuclear use, depending on the distribution of voters who support it. 

Lastly, I include measures of how moral respondents believe the attack to be. One of the questions asks directly whether respondents believe that the nuclear attack was ethical. Another asks whether they think it violates \textit{ahimsa}. These questions are useful in so far as they may reveal consequentialist reasoning. If a respondent believes that an attack is not moral but approves of it nonetheless, then that could imply they are willing to justify the attack based on the military and strategic benefits that it produces.

All outcome variables are measured on seven point Likert scales, with respondents being given the option to select “don't know.” If I were to find evidence to substantiate my hypotheses, I could observe two different, potentially compatible, patterns in people's responses. On one hand, support for nuclear use could decline across the seven point Likert scale. In this scenario, people's support for nuclear use changes with little difference in the number of survey takers reporting that they don't know how to respond. In another scenario, the incidence of people reporting that they don't know how to respond could increase while the number of people supporting the attack decreases. In this case, identity based arguments convert people into fence sitters. Instead of being in favor of military action, they are unsure of how to respond. To account for this possibility, I run all of my analyses twice. In one case, I exclude all NA observations; in the other, I impute values for people who say “don't know,” placing them in the middle of the Likert scale for each of the my outcomes.

Using inverse covariant weighting (ICW), I create an index that represents respondents' cumulative opposition to nuclear use. To create this index, I use the public's approval of the attack, whether it would affect how they vote, and whether they would publicly protest it. I group these factors together because they are indicative of how strongly a politician would be sanctioned for nuclear use. Creating such an index is useful to my study both practically and theoretically. ICW is useful practically because it enhances statistical power and reduces the number of tests that I run, minimizing the multiple comparison problem. Theoretically, ICW is useful because an election-seeking politician will weigh several aspects of public opposition to a nuclear attack when deciding whether to conduct one—not strictly, for example, public disapproval of the attack or whether people would protest it. As a result, though ICW does not provide an easily interpretable outcome variable, the index is holistic and provides insights into the different possible sources of opposition to nuclear use. As such, I estimate the following OLS model as the primary means to evaluate hypotheses two through four:

\begin{equation}
Y_i=\alpha+\beta_1Identity_i+\beta_2NFU_i+\beta_3(Identity_i*NFU_i)+\mathbf{\delta'X_i} +\epsilon_{i},
\end{equation}

\noindent
where $i$ indexes respondents, $Y_i$ represents composite attitudes toward nuclear use, $\beta_1$ is the fixed effect of reading identity-based arguments, $\beta_2$ is the fixed effect of telling respondents about India's NFU, $\beta_3$ is the fixed effect of the interaction between both treatments, and $X_i$ is a vector of individual-level covariates selected using the Belloni, Chernozhukov and Hansen (2014) double-LASSO selection procedure. $Y_i$ is coded such that the higher values reflect stronger opposition to nuclear use.

Although the above model is useful for gauging whether my treatments have an effect on public opposition to nuclear use in the manner that I expect, it has several limitations. Most notably, it cannot be used to evaluate my first hypothesis. To determine whether the public has a categorical aversion to nuclear use, we must look at the sheer magnitude of support in the control group. I do so by collapsing public approval for nuclear use into three categories: disapproval, indifference, and approval. I reduce approval of the attack into three categories instead of two because indifference does not complement arguments either that there is a taboo against nuclear use or that people are particularly receptive to it. As a result, to generate a descriptively accurate assessment of public opposition to nuclear use, I do not collapse indifference into either approval or disapproval of the attack.

Second, even though ICW is a useful way to create a composite score for support of nuclear use, it weights different measures in a manner that's theoretically arbitrary. To be sure, factor analysis or another means of creating a composite score for public opposition to nuclear use would struggle from the same problem. Because different leaders would not weigh different sources of public opposition to nuclear use in the same manner, any particular weighting system that researchers impose would not universally apply. That is to say nothing of the fact that because ICW and factor analysis do not produce easily interpretable results, they would not be useful for political elites deciding whether to use nuclear weapons. 

I account for these problems in two ways. First, to test the robustness of my findings, I implement principal component analysis (PCA). Although this approach does not resolve issues regarding the interpretability of my outcome variable, it will show whether my treatments are still statistically significant using a different weighting scheme, reducing concerns that any statistically significant finding is attributable to one configuration of my composite outcome. Second, I estimate a series of logistic regression models, in which I dichotomize every outcome variable such that indifference to and disapproval of the attack are grouped together. Logistic regression is useful because ICW and PCA are “black boxes.” That is to say, it is not clear which measures of my outcome are being up- or down-weighted. By breaking apart respondents' answers across each of my outcomes, I evaluate how my treatments affect each type of opposition distinct from one another. Although this approach raises the problem of multiple comparisons, it reduces concerns regarding the theoretical arbitrariness of using any given weighting scheme. As such, in addition to my first model, I also estimate the following equation, using each of my measures of my outcome in place of $Y_i$:

\begin{equation}
\log(\frac{\Pr(Y_i=Support)}{\Pr(Y_i\neq Support)})=\alpha+\beta_{1}Arguments_i+\beta_{2}NFU_{i}+\beta_{3}(Arguments_i*NFU_i)+\mathbf{\delta'X_i}
\end{equation}

\noindent
All parameters aside from $Y_i$ represent the same as they did in the previous equation. In text, I report the results for equations (1) and (2). In the Appendix, I include further robustness tests, which I describe in the following section. 

To conduct my experiment, I use Amazon Mechanical Turk to recruit two thousand participants to fill out a survey on Qualtrics. I select MTurk primarily for its cost efficiency compared to other survey platforms. Being cheap to recruit respondents, MTurk enables me to gain enough statistical power to adequately test the hypotheses that I proposed above. That being said, though highly correlated with nationally representative surveys, MTurk's results are not as representative as either probability sampling or online survey panels \citep{berinsky_evaluating_2012}. In India, MTurk's samples are disproportionately young, male, upper-caste, and wealthy \citep{boas_recruiting_2020}. To account for demographic differences, I weight my survey results using entropy balancing. 

That being said, because my survey is conducted in English, it cannot reasonably be construed as representative even after weighting. Because English is most widely spoken among people who are well educated, my findings will likely only apply narrowly to upper-class, elite opinion. In terms of the treatment effect, it is plausible that identity-based appeals could be salient in elite opinion but not in the general population. Given how well educated English-speaking people are on balance, then it is reasonable to believe that they know more about Indian history and foreign policy than the median voter. As a result, identity-based appeals could be more salient in my sample than they would be in the population writ large.

\section{Robustness Tests}
\subsection{Difference of Proportions, Confidence Intervals}
First, instead of using ICW and PCA, I evaluate whether my treatments have their intended effect using difference of proportions, calculating confidence intervals for the point estimates of the amount of support and opposition on each measure of my outcome variable. I do so to get top-line estimates for how my treatments affect overall public support for nuclear use. To conduct the difference of proportions, I use the same dichotomization scheme as described for the logistic regression models. That is to say, disapproval of and indifference to the attack are grouped together, compared to active approval.

\subsection{Different Combinations of Controls}
First, I replicate Equations (1) and (2) using a different combination of controls. First, I ignore all covariates and estimate the effects of my treatments irrespective of them. Subsequently, I remove whatever covariates are not statistically significant using a Type III test and re-estimate treatment effects ignoring them. Third, I include all covariates and re-estimate treatment effects.

\subsection{Partisan Preferences and Gender}
Considering that foreign policy preferences are known to be heavily influenced by respondents' partisan identification and gender, I re-estimate Equation (1), conditioning the effect of the treatments on what party respondents would prefer to lead the Lok Sabha. Likewise, considering that women have generally been found to be more pacifistic than men, I do the same for gender.

\subsection{Dichotomize on Each Step of Outcome Variables}
For each of my outcome variables, I replicate the logistic regression model described by Equation (2) on each one-step increase on the Likert scale. That is to say, I test whether the treatments change the number of people selecting 1 on a given measure of my outcome, compared to 2-7. Subsequently, I test the number of people selecting 1 or 2 on a given measure of my outcome, compared to 3-7. I do so to evaluate whether my treatments are only salient for people with a given strength of belief about the legitimacy of nuclear use.

\bibliography{references}

\end{document}